\section{Preliminaries}
The basic technologies and conceptual structures that underpin contemporary machine-generated medical report production must be understood in order to fully appreciate the proposed solution.

\subsection{Deep Learning in Medical Vision}
Convolutional Neural Networks (CNNs) have been employed for decades in medical image 
processing for obtaining information from MRI and X-ray scans.But the primary focus of new 
developments has been redirected to Vision Transformers (ViT).The transformers employ a 
"self-attention" process,allowing the simulator to figure out the importance of many different 
elements of an image internationally when compared to analyzing very small pixels individually. 
In the field of radiography, where an insignificant irregularity in a particular region of an X-ray 
can correspond with anomalies in another,this capacity is vital. 
\subsection{Large Language Models (LLMs) and the "Black Box" Problem }
Large Language Models (LLMs) are probabilistic mathematical models that were originally 
developed to accurately predict a particular sequence's subsequent token Singhal et al\cite{Singhal2022}. 
highlight the fact that although they are equipped with an extensive amount of knowledge that is 
implicit,that they essentially have no "grounding." An LLM operates as a "Black Box" in a 
healthcare environment;it can come up with a diagnosis that appears credible but does not 
originate from actual information about patients,instead comes from computational 
hallucinations. 

\subsection{Retrieval-Augmented Generation (RAG)} Lewis et al.\cite{Lewis2020}developed the Retrieval-Augmented Generation (RAG) framework to address the grounding problem.In theory,RAG divides the inference procedure into two separate phases:\begin{enumerate} \item \textbf{The Retriever:} In order to find related data regarding an individual's request,an extensive visual search engine searches a database of trustworthy medical records.\item \textbf{The Generator:} The response has been generated using the documents that were obtained as a "set of circumstances" by a sequence-to-sequence framework (the LLM).This segmentation provides assurance that the result generated by the model is a compilation of retrieved information as opposed to solely an educated prediction.\end{enumerate}
\subsection{ Semantic Search and Vector Embeddings}
In the healthcare sector, a traditional keyword search method (BM25) usually does not work because the patients' explain their symptoms in a way that isn't found in textbooks (e.g., "trouble breathing" vs. "dyspnea").In accordance with Kim \cite{Kim2025},Semantic Search uses Vector Embeddings to deal with this. The collection of information and the healthcare search query are simultaneously transformed into multidimensional vector information by the algorithm.The framework is able to retrieve documents that are compatible with the literal meaning of the search query instead of simply the precise phrases through determining the gap between the two via a mathematical procedure called Cosine Similarity.
\section{Introduction and Overview}
From simple organizational tasks to complex report generation tasks, medical diagnosis in the automated sector is developing rapidly.CNN and Transformer,two deep learning methods have shown outstanding results in analyzing medical images.To highlight this development which is clearly unavoidable Sharma et al.\cite{Sharma2020}.AI now not only can detect but also provides a comprehensible description. 
However,some dysfunctionalities are still out there due to relying solely on LLM models. LLM models function in such a way that they are typically known as black box because of their nature of not being able to differentiate between logic and facts. Singhal et al.\cite{Singhal2022} have shown that even though algorithms like Med-PaLM function exceptionally well on standardized assessments,they can be vulnerable to "hallucinations" which end up in trustworthy but factually incorrect medical recommendations.In a highly competitive discipline such as radiology, where just one mistake might result in a life-threatening misdiagnosis,such an absence of factual grounding is beyond unacceptable.
Lewis et al.\cite{Lewis2020}launched Retrieval-Augmented Generation (RAG),a change in perspective that breaks down "memory" (training data) from "reasoning," in order to solve this.In comparison to conventional models that exclusively depend entirely on previously trained stored information, the RAG algorithm facilitates frameworks to retrieve unpredictable evidence from outside.Abo El-Enen et al.\cite{AboElEnen2025} argue in an in-depth investigation of the healthcare RAG that such a structural gap transforms the artificial intelligence coming from a simple "memorizer" into a "reasoner" which bases its decisions on data that can be verified.
\section{Existing Solutions and Technical Limitations
}
According to recent research,medical report generation systems fall into three different generations,each with unique technical shortcomings.
\subsection{Standalone Generative Models}
The initial wave made an effort to train "End-to-End" models.Although the text generated by these systems was flexible,it lacked clinical accuracy. Even sophisticated models like GPT-4 struggle with specific biomedical queries in the absence of external retrieval,as Xiong et al.\cite{Xiong2024}  showed.Additionally,Jeong et al.\cite{Jeong2024} discovered that generative models have a tendency to veer into fabrication in the absence of a "self-reflection" mechanism, in which the model verifies its own work against retrieval data.
\subsection{Naive RAG and Retrieval Noise}
Simple keyword matching was frequently used in early RAG implementations (Naive RAG).A major part of “Retrieval Noise” pointed out by Liang et al.\cite{Liang2025} which is actually a situation where a model is confused between documents which are irrelevant.Ultimately causing incorrect results known as “grounded hallucinations”, Miura et al.\cite{Miura2021}.
\subsection{ The "Biomedical RAG Trilemma"}
Biomedical RAG Trilemma is a formal categorized survey which shows the difficulties during maximizing efficiency while managing three points  (reasoning depth, inference latency,resource constraint) at a time.He et al. \cite{He2025} and Zhao et al\cite{Zhao2026} .
Such as MedRAG, a system that works using a knowledge graph. Yet it follows accuracy over fast and compact algorithms.So rendering them in the real world is not suitable because the environment is already limited in resources.

\section{Bridging the Infrastructure Gap: Resource-Efficient AI Strategies}
In the context of widespread medical artificial intelligence deployment,it has revealed a significant challenge known as “infrastructure gap”.Although modern diagnostic techniques are exhibiting significant potential,they demand extensive computational resources ,particularly GPU-based clusters,which are hardly available at rural health care centers in developing nations. To overcome this limitation, the present work is adopting a “Resource Efficient” design philosophy that emphasizes both diagnostic accuracy and  technological sustainability.This balanced framework aims to ensure the deployment of sophisticated AI-driven diagnostic tools in underdeveloped areas and resource limited environments.

\subsection{Domain Specificity Over Model Scaling
}
Expanding the size of models for improved efficiency has become the most prevalent development in machine learning artificial intelligence (AI). In the field of biomedical science, Gu et al. \cite{Gu2021} challenge this theoretical framework, yet. Their particular domain preliminary training research demonstrates how smaller mathematical models might perform considerably more effectively than vast broadly applicable algorithms when constructed only on healthcare sources (like PubMed). We implement the technique recommended by Gu et al. \cite{Gu2021} to use lightweight, domain-focused architectures rather than a computationally costly generalist model.This guarantees that clinical precision is maintained without compromising system functionality on the small web servers found in rural Thana Health Complexes.

\subsection{Optimizing for Low-Resource Environments}
Clear negotiations must be chosen when creating nations that are developing. Inference Latency and Resource Efficiency receive priority in this framework, which has been guided by the survey findings obtained by He et al.\cite{He2025}  We maintain a small system dimension while gaining access to a large, updateable knowledge base by shifting the "memorization" task from the model's internal parameters to an external retrieval database.
\subsection{Sustainable Semantic Retrieval}
To optimize semantic search in a generative model, matrix extraction is a very good option. As it is a method which balances language and speed.\cite{Kim2025} It opts out the standard keyword search which is insufficient in medical queries.

\section{ Social Trust and Explainable AI (XAI)
}
This research highlights ``Clinical Trust'' with precision. \begin{itemize} \item An area with low resources where people are already living a miserable life can suffer a lot from faulty diagnosis as it is not in their budget to fix errors other than their own. So it led to a critical situation for both doctor and patient, as patients are not able to pay for the faulty treatment. Miura et al.\cite{Miura2021}. \item As we stated  earlier about black box situations,because of this doctors cannot trust completely in AI, Rosendbacke et al.\cite{Rosenbacke2024} So, Vries et al \cite{DeVries2023} proposed an idea about implementing RAG which will help in generating an image report which verifies with previous records allowing a doctor to cross check with evidence based reports.
 \end{itemize}

\section{Summary of Key Findings}
The following points represent the present situation of medical diagnostics in the automated sector: \begin{enumerate} \item \textbf{The superiority of RAG:} In healthcare environments, Retrieval-Augmented Generation has effectively substituted independently operated LLMs. RAG significantly decreases the frequency of "hallucinations" when compared with traditional End-to-End frameworks by distinguishing logic from memory \cite{Lewis2020,AboElEnen2025}.  \item \textbf{The Accuracy Latency trade-off:} The trade-off between accuracy and latency in modern frameworks is commonly referred to as the ``Biomedical RAG Trilemma.'' Significant precision for diagnosis systems can frequently lag and demand resources because of the high computational latency \cite{He2025,Zhao2026}. \item \textbf{Possibility of Domain Specificity:} Studies have demonstrated that, after adequate prior training, more compact domain-specific algorithms (such as PubMedBERT) could potentially be able to produce results comparable to more complex models\cite{Gu2021} . \end{enumerate}

\section{ Research Gap}
The scientific research on "Lightweight, Web-Based Architectures" for environments with limited resources continues to have a significant gap in spite of the above findings.
The research studies which have been previously conducted,like those by Zhao et al.  \cite{Zhao2025} and Xiong et al.\cite{Xiong2024}, predominantly focus on theoretical benchmarks with advanced GPU clusters. Investigation on carrying out these types of designs with the limited computer hardware that exists in nations that are developing is inadequate. The "Infrastructure Gap", the challenge in deploying cutting-edge diagnostic artificial intelligence on traditional web hosting platforms while maintaining healthcare precision, has not been sufficiently addressed by the research currently in publication.
\section{Conclusion}
In conclusion,there are several solid theoretical foundations for healthcare artificial intelligence (AI), yet there are not any practical, limited-resource applications for it.Even though the transition from "detection" to "description" has already been implemented in circumstances with plenty of resources, countries with limited resources continue to be incapable of taking advantage of it. For the purpose to offer readily available diagnostics for rural areas of Bangladesh, this thesis constructs a resource-friendly RAG framework based on the specific to the domain strategies recommended by Gu et al.\cite{Gu2021}.
